\paragraph{地形模型的输入与拟合机制。}
本分支的$X_t$使用预测起点锚定的东西／南北坐标系。
令$Z_t$为至多三个可见或生成位置及其时间的缓冲。
首末割线给出速度$v_t$，归一化得到历史方向$h_t$，$\eta_t$表示方向是否有效。
历史设计为$r(Z_t)=(1,v_{t,e},v_{t,n},\|v_t\|,h_{t,e},h_{t,n},\eta_t)^\top$。
地形漂移及位置动力学为
\begin{equation}\label{eq:main-terrain-model}
 b_g(X_t,Z_t)=B_0^\top r(Z_t)+\theta_g^\top\widetilde\phi_g(X_t,Z_t),
 \qquad dX_t=b_g(X_t,Z_t)\,dt+L_g\,dW_t .
\end{equation}
$\widetilde\phi_g$是在选定地图／历史特征后追加截距的向量，包含编码数值及有效性掩码。
道路和河流输入描述最近要素距离与方向，覆盖采用分组指示，坡面输入描述坡度／方向，
登记交互组合部分地图和历史量。这些是固定工程表示；系数从位移率学习，
不受道路吸引或河流回避等物理定律约束。
\paragraph{实际选入的因子，而不是笼统的“水、路”。}
道路组（road）来自保留的Overture道路线；该道路查询排除
bridleway、cycleway、footway、path、pedestrian、steps、track类别，
所以它不等于完整步行路网。河流组（river）来自HydroRIVERS河流线，
不是所有水面、湖岸、可渡性或桥梁通行信息。地表覆盖组（cover，登记键worldcover）
来自WorldCover像素；坡面几何组（surface）来自保留DEM规则的坡度、坡向与法向。
绝对高程只作为计算地表几何的来源，不是本次选入的直接条件器列；
其他路径、铁路、湖岸距离、曲率和粗糙度表示未进入这十个配置。

\paragraph{距离、方向和历史量具体如何组合。}
在当前查询位置$X$，令$d_g$为到最近登记线要素点的距离（米），
$u_g=(u_{g,e},u_{g,n})$为从查询位置指向该点的单位向量，$g\in\{\mathrm{road},\mathrm{river}\}$。
最近点在提供器的局部米制投影中选择，报告距离按其球面距离规则计算。
输入既保留两个方向分量，也保留对数距离及距离--方向乘积：
\begin{equation}\label{eq:main-map-distance}
 \ell_g=\log(1+d_g/(1\,\mathrm m)),\qquad c_g=\ell_g u_g .
\end{equation}
$u_g$不是沿线切向量；对河流也不是流向。$c_g$是两个数值乘积列，
不是道路吸引力或河流排斥力的硬编码。
本次没有$u_{\rm river}^{\top}h$、$u_{\rm river}^{\top}v$或河流距离乘历史速度的直接输入。
历史速度仍通过共享漂移进入预测；地图与运动还会通过生成位置的反馈间接耦合。

地表覆盖向量$w$是四组指示：建成区50；植被10、20、30、40、95、100；
水／湿地80、90；裸地／雪／其他有效有限代码。无数据不算有效“其他”。
坡面向量为$s=(\beta,\sin\alpha,\cos\alpha,n_e,n_n,n_u)$，
$\beta$为弧度坡度，$\alpha$从北顺时针表示下坡坡向，$n$为三维向上的单位法向。
令上坡方向$a=-(\sin\alpha,\cos\alpha)$，选定等高线方向$q=(-a_n,a_e)$。
四项登记交互准确为
\begin{equation}\label{eq:main-map-compositions}
 c_{\rm road}=\ell_{\rm road}u_{\rm road},\quad
 c_{\rm river}=\ell_{\rm river}u_{\rm river},\quad
 c_{\rm cover}=\ell_{\rm road}w,\quad
 c_{\rm surface}=(h^\top a,h^\top q).
\end{equation}
最后一项才直接组合地形与历史方向；它不再乘坡度或速度大小，也不是六乘二的全部叉积。
反向选择$q$会改变第二个投影的符号；平地坡向记录为零，不代表真实的垂直运动关系。

\begin{table}[!htbp]\centering\small
\caption{地形分支的具名因子与实际数值输入。维度不含有效性掩码；同一数值列另有一个有效性指示。
基础列共18维，登记交互共10维，全部配置共28维数值。}
\label{tab:main-terrain-factors}
\begin{tabular}{@{}p{.17\linewidth}p{.43\linewidth}p{.32\linewidth}@{}}\toprule
组名／来源 & 基础因子（数值维度） & 依赖该组的交互\\\midrule
history／因果历史 & 东／北单位历史方向$h$（2） & 与surface共同组成$(h^\top a,h^\top q)$（2）\\
road／道路 & 对数距离$\ell_{\rm road}$（1）、指向最近道路点$u_{\rm road}$（2） & $\ell_{\rm road}u_{\rm road}$（2）；与cover共同组成$\ell_{\rm road}w$（4）\\
river／河流 & 对数距离$\ell_{\rm river}$（1）、指向最近河流点$u_{\rm river}$（2） & $\ell_{\rm river}u_{\rm river}$（2）；没有直接历史速度乘积\\
cover／地表覆盖 & 四组WorldCover指示$w$（4） & 与road共同组成$\ell_{\rm road}w$（4）\\
surface／坡面几何 & 坡度、下坡坡向两分量、法向三分量$s$（6） & 与history共同组成两个有符号方向投影（2）\\\bottomrule
\end{tabular}\end{table}

\paragraph{base、all、LOO与LIO不是额外因子。}
记$B$为base历史配置，$A$为all-terrain全组配置。
LOO（leave-one-group-out，移除一组）是$A-g$；
LIO（leave-one-group-in，仅加入一组）是$B+g$，并非不含历史的单因素模型。
表\ref{tab:main-terrain-configurations}连接这些操作与结果表的简写。
所有配置保留history；删除一组时，同时删除依赖它的交互，然后重新拟合条件器和扩散。

\begin{table}[!htbp]\centering\small
\caption{十配置与实际保留因子／交互的对应。$K$为条件器数值列数，输入宽度为$2K$（数值及掩码），
另加一个拟合截距；独立的七列历史漂移不计入$K$。结果表与图采用$B$、$A-g$、$B+g$，原执行键在审查记录保留。}
\label{tab:main-terrain-configurations}
\begin{tabular}{@{}p{.23\linewidth}p{.28\linewidth}p{.35\linewidth}r@{}}\toprule
配置／操作 & 保留地图组（均有history） & 保留交互 & $K$\\\midrule
$B$／base & 无 & 无 & 2\\
$A$／all & road、river、cover、surface & 全部四项 & 28\\
$A-\mathrm{road}$／LOO & river、cover、surface & 河流距离--方向；坡面--历史 & 19\\
$A-\mathrm{river}$／LOO & road、cover、surface & 道路距离--方向；覆盖--道路；坡面--历史 & 23\\
$A-\mathrm{cover}$／LOO & road、river、surface & 道路及河流距离--方向；坡面--历史 & 20\\
$A-\mathrm{surface}$／LOO & road、river、cover & 道路及河流距离--方向；覆盖--道路 & 20\\
$B+\mathrm{road}$／LIO & road & 道路距离--方向 & 7\\
$B+\mathrm{river}$／LIO & river & 河流距离--方向 & 7\\
$B+\mathrm{cover}$／LIO & cover & 无（不含road） & 6\\
$B+\mathrm{surface}$／LIO & surface & 坡面--历史（history仍在） & 10\\\bottomrule
\end{tabular}\end{table}
覆盖--道路交互在移除road或cover时都会被移除；两个删除损失不能相加成独立因素贡献。
仅cover配置没有道路距离，因此也没有该交互。LOO与LIO考察不同背景下重新拟合的预测配置，
不是同一已拟合模型上把某一输入置零，也不能分别识别距离、方向或某个覆盖子类的因果效应。
\FloatBarrier


共享$B_0$先以强度$10^{-6}$的ridge拟合。
404个外层训练窗口各贡献一个实际首未来转移$(\Delta x_i,\Delta t_i)$，
经过时长不必等于方法观测间隔$\tau$。
随后配置$g$对剩余位移率$y_i=\Delta x_i/\Delta t_i-B_0^\top r(Z_i)$拟合
\begin{equation}\label{eq:main-terrain-fit}
 \widehat\theta_g=\underset{\theta}{\operatorname{argmin}}
 \left\{\frac1{2n}\|D_g\theta-Y\|_F^2+
             \frac\lambda2\|\theta\|_F^2\right\},
 \quad n=404,\quad\lambda=10^{-4} .
\end{equation}
$D_g$堆叠增广特征行。两阶段拟合对参数线性，但地图查询和历史反馈使位置动力学依赖状态。
各配置分别从拟合位移残差估计扩散：
$Q_g=L_gL_g^\top=n^{-1}\sum_i
(\Delta x_i-\widehat b_{g,i}\Delta t_i)
(\Delta x_i-\widehat b_{g,i}\Delta t_i)^\top/\Delta t_i$。
这与方法分支的中心化$\tau R$估计不同。
Euler--Maruyama逐路径加入漂移$b_g\Delta$与噪声$\sqrt\Delta L_gz$；
每5个保存秒向历史缓冲追加生成位置并重算割线，地图在预测位置查询。
真实未来观测不会进入缓冲。

base配置保留共享历史漂移及拟合的纯历史残差条件器，不查询地图图层；
all-terrain加入全部登记组及交互，LOO从该完整配置移除一组，
LIO向历史基线加入一组。因此RQ2比较这些配置在匹配目标上的预测增量。
共享基线与分开的残差拟合明确了所考察的增量，重拟合和特征惩罚则限制单个因素归因。
